\documentclass[11pt]{article}

\usepackage[T1]{fontenc}
\usepackage[utf8]{inputenc}
\usepackage{lmodern}
\usepackage{microtype}
\usepackage[a4paper,margin=25mm,headheight=15pt]{geometry}
\usepackage{amsmath,amssymb,mathtools}
\usepackage{booktabs,tabularx,longtable,array}
\usepackage{enumitem}
\usepackage{xcolor}
\usepackage{hyperref}
\usepackage{fancyhdr}
\usepackage{titlesec}
\usepackage{listings}
\usepackage{tikz}
\usetikzlibrary{arrows.meta,positioning,fit}

\definecolor{navy}{HTML}{15395B}
\definecolor{steel}{HTML}{4A6682}
\definecolor{pale}{HTML}{ECF3F8}
\definecolor{green}{HTML}{25633E}
\definecolor{greenpale}{HTML}{EAF5EE}
\definecolor{amber}{HTML}{7A5600}
\definecolor{amberpale}{HTML}{FFF4D6}
\definecolor{red}{HTML}{8D2B35}
\definecolor{redpale}{HTML}{FBECEE}
\definecolor{ink}{HTML}{202C38}

\hypersetup{
  colorlinks=true,
  linkcolor=navy,
  urlcolor=steel,
  citecolor=navy,
  pdftitle={Finishing the forall-PO-Free Fragment of ALCI-RCC5},
  pdfsubject={Fresh-occurrence cone-scheme decision procedure and Lean certification blueprint}
}

\pagestyle{fancy}
\fancyhf{}
\fancyhead[L]{\small\color{steel}$\forall\PO$-free $\mathcal{ALCI}_{\mathrm{RCC5}}$}
\fancyhead[R]{\small\color{steel}Certification plan}
\fancyfoot[C]{\small\thepage}
\renewcommand{\headrulewidth}{0.4pt}

\titleformat{\section}{\Large\bfseries\color{navy}}{\thesection}{0.65em}{}
\titleformat{\subsection}{\large\bfseries\color{navy}}{\thesubsection}{0.65em}{}
\titleformat{\subsubsection}{\normalsize\bfseries\color{steel}}{\thesubsubsection}{0.65em}{}
\setlist[itemize]{leftmargin=1.6em,itemsep=0.25em,topsep=0.35em}
\setlist[enumerate]{leftmargin=1.8em,itemsep=0.25em,topsep=0.35em}
\setlength{\parindent}{0pt}
\setlength{\parskip}{0.55em}
\allowdisplaybreaks

\lstdefinelanguage{Lean}{
  morekeywords={theorem,def,structure,inductive,where,by,let,have,show,from,
    if,then,else,match,with,fun,forall,exists,namespace,end,import,open,
    instance,deriving,Decidable,Finite,Fintype,Prop,Type},
  sensitive=true,
  morecomment=[l]{--},
  morestring=[b]"
}
\lstset{
  language=Lean,
  basicstyle=\ttfamily\small,
  keywordstyle=\color{navy}\bfseries,
  commentstyle=\color{steel},
  stringstyle=\color{green},
  columns=fullflexible,
  keepspaces=true,
  breaklines=true,
  frame=single,
  rulecolor=\color{steel!40},
  backgroundcolor=\color{pale!55},
  xleftmargin=0.4em,
  xrightmargin=0.4em,
  literate={∈}{{$\in$}}1
}

\newcommand{\PP}{\mathsf{PP}}
\newcommand{\PPI}{\mathsf{PPI}}
\newcommand{\PO}{\mathsf{PO}}
\newcommand{\DR}{\mathsf{DR}}
\newcommand{\EQ}{\mathsf{EQ}}
\newcommand{\Ty}{\mathsf{Ty}}
\newcommand{\Cl}{\mathsf{Cl}^{+}}
\newcommand{\cone}{\operatorname{cone}}
\newcommand{\At}{\mathsf{At}}
\newcommand{\Sat}{\mathsf{Sat}}
\newcommand{\Phiop}{\Phi}
\newcommand{\statusbox}[3]{%
  \begin{center}
  \fcolorbox{#1}{#2}{\begin{minipage}{0.91\linewidth}\textbf{#3}\end{minipage}}
  \end{center}}

\begin{document}

\pagenumbering{gobble}
\begin{titlepage}
\thispagestyle{empty}
\vspace*{1.5cm}
{\color{navy}\rule{\linewidth}{1.2pt}}\\[1.1cm]
{\huge\bfseries\color{navy}Finishing the $\forall\PO$-Free Fragment\\[0.22cm]
of $\mathcal{ALCI}_{\mathrm{RCC5}}$\par}
\vspace{0.7cm}
{\Large\color{steel}A fresh-occurrence cone-scheme proof strategy\\
and Lean certification blueprint\par}
\vspace{1.35cm}

\begin{center}
\renewcommand{\arraystretch}{1.25}
\begin{tabularx}{0.91\linewidth}{>{\bfseries}rX}
Purpose: & detailed replacement plan for a complete decision proof \\
Input reviewed: & round-two cold-attack package, 2026-08-28 \\
Report date: & 2026-08-28 \\
Formal status: & proof architecture and machine-checkable blueprint; Lean proof pending
\end{tabularx}
\end{center}

\vfill
\noindent\fcolorbox{navy}{pale}{%
\begin{minipage}{0.93\dimexpr\linewidth-2\fboxsep-2\fboxrule\relax}
\textbf{Recommendation.}
Retire the finite-kernel and borrowed-edge extraction target. Decide the fragment by a
finite greatest-fixed-point computation over \emph{cone signatures}, then unfold a
surviving finite strategy into an infinite fresh-occurrence forest. Proper-part edges are
created only at birth; disconnection is the downward closure of designated horizontal
birth cuts; overlap is residual and protected at its birth cut. This removes the three
known obstructions by construction and yields a direct soundness/completeness theorem.
\end{minipage}}\par
\vspace{0.95cm}
{\color{navy}\rule{\linewidth}{1.2pt}}
\end{titlepage}

\pagenumbering{roman}
\begin{abstract}
This report gives a concrete route to a full Lean certification of decidability for the
fragment of $\mathcal{ALCI}_{\mathrm{RCC5}}$ in negation normal form whose closure contains
no universal $\PO$ restriction. The route deliberately abandons finite quotient models,
shared witness pools, and borrowed edges. Its finite object is instead a regular
\emph{strategy}: a cone signature consists of a local support type and a finite spectrum
of possible strict-lower types. A monotone elimination operator retains precisely those
signatures whose existential demands have compatible target signatures. The survivor set
is finite and computable.

Soundness unfolds the strategy into fresh occurrences. Vertical birth edges generate a
ranked strict order; $\DR$ births generate separated vertical components; $\PO$ births
are protected as residual non-order, non-disjoint bridges. A one-way truth lemma realizes
every selected type. Completeness maps any source model to its realized cone signatures:
RCC5 singleton composition cells force the required $\PP$, $\PPI$, and $\DR$ transfer
conditions, while the syntactic exclusion of $\forall\PO$ removes every $\PO$ transfer
obligation. Realized signatures form a post-fixed point and hence survive elimination.

The result is a finite decision procedure with a conservative double-exponential upper
bound. The package accompanying this report contains the state rules, a Lean module map,
a proof-obligation ledger, and executable structural checks. It is a formal certification
plan, not a claim that the final Lean theorem already compiles.
\end{abstract}

\setcounter{tocdepth}{1}
\tableofcontents
\clearpage
\pagenumbering{arabic}

\section{Executive answer}

\statusbox{green}{greenpale}{No counterexample was found to the replacement architecture.
It is a complete candidate soundness/completeness route: the finite state space and
acceptance test are explicit, and every difficult obligation has a named Lean theorem.
Occurrence geometry and the truth lemma remain genuine go/no-go research lemmas until
they compile.}

The final public theorem should state:
\[
  \forall C,\quad \mathsf{ForallPOFree}(C)\longrightarrow
  \mathsf{Decidable}\bigl(\exists M,x,\;M,x\models C\bigr).
\]
It should be proved through a Boolean core
\[
  \mathsf{coneSatCoreB}(C)=\mathsf{true}
  \quad\Longleftrightarrow\quad
  \exists M,x,\;M,x\models C.
\]
under the displayed fragment hypothesis. A separate three-valued wrapper should report
\texttt{outOfScope}, rather than \texttt{unsat}, when that hypothesis is absent.

The decisive change is to separate \emph{finite control} from \emph{model domain}.
There need not be a finite model. A finite control state may recur indefinitely, but every
existential use creates a fresh domain occurrence. Therefore a $\PP$ self-transition in
the finite control graph denotes an infinite ascending chain, never a reflexive
proper-part edge.

\begin{center}
\renewcommand{\arraystretch}{1.25}
\begin{tabularx}{\linewidth}{>{\bfseries}p{0.24\linewidth}XX}
\toprule
Issue in the previous route & Replacement invariant & Why it closes the issue \\
\midrule
D1: cycles from borrowed order edges &
Only fresh parent--child vertical birth edges; a height map increases on every base edge. &
The transitive closure is irreflexive by an integer-rank argument. \\
D2: blocked source and target need incompatible order directions &
A demand is realized by a new occurrence even when its target control state repeats. &
Control-state equality never identifies domain points or reverses an old edge. \\
E: anchors create crossed $\PO/\PP/\DR$ constraints &
No finite external anchors or globally shared witnesses exist. &
Each recursive demand unfolds below its own occurrence; horizontal cuts are local. \\
\bottomrule
\end{tabularx}
\end{center}

\section{Scope and exact semantic target}

\subsection{Syntax}

The public wrapper accepts a raw concept and first proves semantic normalization:
\[
  M,x\models C\quad\Longleftrightarrow\quad M,x\models\mathsf{nnf}(C).
\]
Inverse-role syntax is normalized at the same boundary by replacing $r^{-}$ with the
RCC5 converse atom. Define $\mathsf{ForallPOFree}(C)$ by inspecting the positive closure
of $\mathsf{nnf}(C)$, not the surface syntax.

Let $\Cl$ be the positive subformula closure of $\mathsf{nnf}(C)$, including every modal
body needed by structural induction. It is deliberately not completed with Boolean duals.
Dual-complete types would force preservation of negative $\PO$ information such as
$\neg\exists\PO.D=\forall\PO.\neg D$, which the residual-$\PO$ construction does not
support. Write
\[
  m=|\Cl|.
\]

The role alphabet is
\[
  \At=\{\EQ,\PP,\PPI,\DR,\PO\}.
\]
The input condition is
\[
  \forall D,\quad \forall\PO.D\notin\Cl.
\]
All other existential and universal role restrictions are permitted. Converse is already
represented by $\PP/\PPI$; $\EQ$ is strong identity.

The theorem concerns \emph{concept satisfiability only}: no TBoxes, ABoxes, nominals,
counting restrictions, role conjunction, or identity-bearing assertions are included.
The only roles after normalization are the five RCC5 atoms above.

\subsection{Frames and models}

The primary semantic interface is an abstract atomic RCC5 network: exactly one atom holds
on each ordered pair, $\EQ$ is equality, converse is respected, and every triangle obeys
the RCC5 composition table. The intended normal form is an ordered-disjoint frame:

\begin{itemize}
  \item a strict partial order $<$ representing $\PP$;
  \item a symmetric, irreflexive disjointness relation $\perp$;
  \item downward closure:
    $x'\leq x$, $y'\leq y$, $x\perp y$ imply $x'\perp y'$;
  \item $\PPI$ as converse, $\EQ$ as identity, and $\PO$ as the residual case.
\end{itemize}

The existing or separately proved normal-form theorem then supplies the entire RCC5
composition table. For a self-contained semantic endpoint, Section~\ref{sec:tokens}
gives an exact representation by arbitrary nonempty sets. It does not establish
regular-closed or topological region representability.

\subsection{What this report does and does not certify}

The report specifies the definitions, intermediate lemmas, theorem dependency order,
finite bounds, regression obligations, and release gates needed for a final Lean
certificate. The included scripts validate consistency of the blueprint and replay the
known finite adversarial instances. They do not replace the missing Lean sources. No
\texttt{sorry}-free theorem is claimed until the module graph in
Section~\ref{sec:lean} compiles and the audit gates in Section~\ref{sec:audit} pass.

\section{Why the target must be a regular strategy}

\subsection{The fragment lacks the finite model property}

Let
\[
  H=\exists\PP.\top,\qquad
  C_{\infty}=H\wedge\forall\PP.H.
\]
This concept is $\forall\PO$-free. It is satisfiable on the strict chain of finite initial
segments
\[
  \{0\}\;\PP\;\{0,1\}\;\PP\;\{0,1,2\}\;\PP\;\cdots.
\]
But no finite strict transitive model can satisfy it. From a root satisfying $H$ choose a
proper superpart; the universal conjunct ensures that each proper superpart again has a
proper superpart. Finiteness would force repetition, contradicting irreflexivity.

Therefore no proof based on a bounded finite quotient domain can be complete. The finite
certificate must represent an infinite regular model, exactly as a finite automaton
represents infinitely many runs.

\subsection{Control recurrence is not point identification}

Suppose a finite state $q$ has a $\PP$ transition to itself. The unfolding has occurrences
\[
  \epsilon,\quad d,\quad dd,\quad ddd,\ldots
\]
all labelled by $q$, with strictly increasing height. This is sound. Identifying them
would instead assert $x\PP x$, which is impossible. The Lean development must enforce
this distinction at the type level: occurrences are finite demand paths, while control
states are labels attached to paths.

\subsection{Why no parity condition is needed}

ALCI has no least or greatest fixpoint operators. Every existential demand is fulfilled
at the next generation, and universal restrictions are safety invariants checked along
the generated relations. Infinite recurrence therefore carries no fairness or eventuality
obligation. The elimination procedure is a one-priority safety game, not a parity game.

\section{Finite cone signatures}

\subsection{Support types}

A support type $T\subseteq\Cl$ is locally Hintikka-consistent:

\begin{itemize}
  \item $\bot\notin T$, and $\top\in T$ whenever $\top\in\Cl$;
  \item complementary concept literals do not both occur;
  \item $D\wedge E\in T$ implies $D,E\in T$;
  \item $D\vee E\in T$ implies $D\in T$ or $E\in T$;
\end{itemize}

Types are support sets, not maximal or Boolean-complete theories. The semantic induction
proves only membership-to-truth; no converse truth lemma is needed.

Let $\Ty$ be the finite set of support types and let
\[
  N=|\Ty|\leq 2^m.
\]
For $r\in\At$ define the universal body set
\[
  A_r(T)=\{D\in\Cl\mid \forall r.D\in T\}.
\]

\subsection{State definition}

A cone signature is a pair
\[
  q=(T,S),\qquad T\in\Ty,\quad S\subseteq\Ty.
\]
Its inclusive cone is
\[
  \cone(q)=S\cup\{T\}.
\]
The intended invariant is that every strict lower occurrence of a point labelled by $q$
has a support type in $S$. The spectrum is an upper bound, not an exact list; this
monotonicity is important in soundness.

The number of raw signatures is
\[
  K=N\,2^N\leq 2^{m+2^m}.
\]

\subsection{Static local admissibility}

A state $q=(T,S)$ is statically admissible when:

\begin{enumerate}
  \item $T$ is a support type;
  \item $\forall\EQ.D\in T$ implies $D\in T$;
  \item $\exists\EQ.D\in T$ implies $D\in T$;
  \item for every $U\in S$,
  \[
    A_{\PP}(U)\subseteq T
    \qquad\text{and}\qquad
    A_{\PPI}(T)\subseteq U.
  \]
\end{enumerate}

The fourth condition is the entire vertical universal check. If a point of type $U$ is a
proper part of a point of type $T$, then $\forall\PP$ bodies at the lower point and
$\forall\PPI$ bodies at the upper point must hold at the other endpoint.

\subsection{Transition compatibility}

Let $\delta=\exists r.D\in T_q$ be a demand at state $q$, and let $q'$ be a proposed target.
First require $D\in T_{q'}$. The role-specific conditions are:

\begin{center}
\renewcommand{\arraystretch}{1.35}
\begin{tabularx}{\linewidth}{>{\bfseries}p{0.11\linewidth}X}
\toprule
Role & Additional compatibility condition \\
\midrule
$\EQ$ & No child is created; require $D\in T_q$. \\
$\PP$ & $\cone(q)\subseteq S_{q'}$. \\
$\PPI$ & $\cone(q')\subseteq S_q$. \\
$\DR$ & For all $U\in\cone(q)$ and $V\in\cone(q')$,
$A_{\DR}(U)\subseteq V$ and $A_{\DR}(V)\subseteq U$. \\
$\PO$ & No additional transfer condition. \\
\bottomrule
\end{tabularx}
\end{center}

The $\PO$ row is not a shortcut. There is no formula $\forall\PO.E$ in the closure, so a
$\PO$ edge carries no universal body. The unfolding must still guarantee that the
designated endpoints are exactly $\PO$; that is a geometric theorem, not a type-transfer
condition.

\section{The finite greatest-fixed-point decider}

Let $\mathsf{Static}$ be the finite set of statically admissible signatures. For
$X\subseteq\mathsf{Static}$ define
\[
 \Phiop(X)=
 \left\{q\in X\ \middle|\ 
 \begin{array}{l}
 \text{for every non-$\EQ$ existential demand $\delta$ in $T_q$,}\\
 \text{there exists $q'\in X$ compatible with $(q,\delta)$, and}\\
 \text{every $\EQ$ demand is locally satisfied}
 \end{array}
 \right\}.
\]
This operator is monotone and reductive. Compute
\[
  X_0=\mathsf{Static},
  \qquad X_{i+1}=\Phiop(X_i).
\]
The sequence descends and stabilizes in at most $|\mathsf{Static}|\leq K$ strict
iterations. Let $X_\infty$ be the stable set.

\textbf{Acceptance criterion.}
The raw input $C_0$ is accepted iff some $q\in X_\infty$ contains the prepared closure
index of $\mathsf{nnf}(C_0)$.

\textbf{Finite certificate.}
For every survivor $q$ and every non-$\EQ$ demand in $T_q$, choose one compatible target
in $X_\infty$. The resulting strategy has at most $K$ states and at most $mK$
chosen transitions. The abstract scheme uses a proof-independent total target function
from state and full role-aware demand code to state; membership and correctness are
separate fields.

\begin{lstlisting}
def coneSatCoreB (C : Concept) : Bool :=
  let cl      := positiveClosure (nnf C)
  let states  := enumerateStaticStates cl
  let survive := iterateAtMost states.card (eliminate cl) states
  survive.any fun q => q.ty.contains (nnf C)

theorem coneSatCoreB_correct (hfree : ForallPOFree C) :
    coneSatCoreB C = true <-> Satisfiable C

def coneCheck (C : Concept) : CheckResult :=
  if forallPOFreeB C then
    if coneSatCoreB C then .sat else .unsat
  else .outOfScope

theorem coneCheck_sat_iff :
    coneCheck C = .sat <->
      ForallPOFree C /\ Satisfiable C
theorem coneCheck_unsat_iff :
    coneCheck C = .unsat <->
      ForallPOFree C /\ Not (Satisfiable C)
theorem coneCheck_outOfScope_iff :
    coneCheck C = .outOfScope <-> Not (ForallPOFree C)
\end{lstlisting}

The proof of \texttt{coneSatCoreB\_correct} is split into the soundness and completeness
constructions below. The executable procedure must use reflected Boolean predicates with
correctness theorems, rather than classical existence tests hidden behind opaque
decidability.

\section{Fresh-occurrence unfolding}

\subsection{Occurrence domain}

Fix a chosen transition function for the survivors. An occurrence is a finite valid word
of demand choices. The root is the empty word. If occurrence $x$ has control state $q$
and $\delta$ is a non-$\EQ$ demand in $T_q$, then $x\mathbin{\raisebox{0.2ex}{\(\smallfrown\)}}\delta$
is a distinct child labelled by the selected target state.

Two occurrences with the same control state remain different domain elements. Each
non-$\EQ$ demand at each occurrence receives its own fresh child. A transition pointer is
therefore only a control pointer; it is never a semantic backpointer.

\subsection{Vertical birth order and rank}

Orient base order edges according to the demand role:

\begin{itemize}
  \item a $\PP$ child lies above its parent: $x<_{0}x\delta$;
  \item a $\PPI$ child lies below its parent: $x\delta<_{0}x$;
  \item $\DR$ and $\PO$ children start new vertical components;
  \item $\EQ$ creates no child.
\end{itemize}

Let $<$ be the transitive closure of $<_{0}$. Define a height
$h:\mathsf{Occurrence}\to\mathbb{Z}$ by adding $1$ at a $\PP$ birth, subtracting $1$ at a
$\PPI$ birth, and resetting the relative origin at a horizontal birth. Every oriented
base edge satisfies
\[
  x<_{0}y\Longrightarrow h(x)+1=h(y).
\]
Thus every nonempty $<_{0}$ path strictly raises height, and $<$ is irreflexive.
Transitivity holds by definition. Parent uniqueness of words gives the induction scheme
needed for this rank proof.

\subsection{Vertical components}

Give each occurrence a component tag equal to the sequence of uniquely identified
horizontal birth edges ($\DR$ or $\PO$, including the child occurrence) on its ancestry.
Vertical births preserve the tag; horizontal births append a fresh edge identifier.
Order edges never cross component tags. Removing horizontal births yields vertical
components, and contracting each such component yields a tree whose edges are the unique
horizontal birth bridges.

This component tree is the geometric skeleton. It replaces every global anchor, witness
pool, and kernel-incidence construction.

\subsection{Disjointness closure}

For each $\DR$ birth bridge $(x,y)$, declare the two inclusive vertical cones at that
bridge disjoint. Equivalently, take the symmetric downward closure of the seed bridges:
\[
  u\perp v
  \quad\Longleftrightarrow\quad
  \exists(x,y)\in\mathsf{DRSeed},\
  (u\leq x\wedge v\leq y)\ \vee\
  (u\leq y\wedge v\leq x).
\]
Here $\leq$ is the reflexive closure of $<$.

The central separator lemma states that the two sides of a $\DR$ birth cut have no common
lower occurrence. It follows by unique component ancestry: any common lower would have
to cross the unique horizontal bridge, but order never crosses component tags. Hence
$\perp$ is irreflexive. Symmetry and downward closure are immediate, and
\[
  x<y\Longrightarrow \neg(x\perp y).
\]

\subsection{Protected overlap}

Define RCC5 atoms from the order and disjointness normal form:
\[
  R(x,y)=
  \begin{cases}
  \EQ & x=y,\\
  \PP & x<y,\\
  \PPI & y<x,\\
  \DR & x\perp y,\\
  \PO & \text{otherwise.}
  \end{cases}
\]

At a $\PO$ birth bridge, the endpoint occurrences differ, lie in different vertical
components, and are not seeded disjoint. Unique-bridge separation prevents any later
vertical path between them; downward $\DR$ closure cannot cross a non-$\DR$ bridge.
Therefore the birth endpoints remain neither ordered nor disjoint and hence are exactly
$\PO$. This is the \emph{protected-PO theorem}.

It is important to prove this theorem from component ancestry before defining $\PO$ as
residual. An incremental default-$\PO$ labelling algorithm would be unsound because later
order or disjointness closure could overwrite an early guess.

\subsection{Exact set representation}\label{sec:tokens}

If the public semantics require literal nonempty sets, represent the ordered-disjoint
frame without quotient types. Let
\[
  \mathsf{Tok}=\{(a,b)\mid \neg(a\perp b)\}
\]
and define
\[
  \mathcal{R}_x
  =\{(a,b)\in\mathsf{Tok}\mid a\leq x\ \vee\ b\leq x\}.
\]
Then:
\[
  \mathcal{R}_x\subseteq\mathcal{R}_y\Longleftrightarrow x\leq y,
  \qquad
  \mathcal{R}_x\cap\mathcal{R}_y=\varnothing\Longleftrightarrow x\perp y.
\]
The self token $(x,x)$ makes $\mathcal{R}_x$ nonempty and distinguishes incomparable
points. Ordered pairs are a useful proof-engineering choice: they avoid quotient
soundness obligations and retain both sides of a non-disjoint witness. These equivalences
transfer the ordered-disjoint atom exactly to standard RCC5 relations on nonempty sets.

\section{Soundness: a survivor yields a model}

\subsection{Spectrum invariant}

Let $\sigma(x)$ be the control state at occurrence $x$. Prove simultaneously with the
unfolding that
\[
  y<x \Longrightarrow T_{\sigma(y)}\in S_{\sigma(x)}.
\]
For one vertical birth this is exactly the $\PP$ or $\PPI$ transition rule. For a
transitive path, the inclusive-cone subset propagates the earlier endpoint type into the
later spectrum. This is why the transition rules use $\cone(q)$ rather than only $S_q$.

\subsection{Universal restrictions}

Assume $D\in T_{\sigma(x)}$ is a universal body obligation.

\begin{itemize}
  \item If $\forall\PP.D\in T_{\sigma(x)}$ and $x<y$, the spectrum invariant and static
        local admissibility put $D$ in $T_{\sigma(y)}$.
  \item If $\forall\PPI.D\in T_{\sigma(x)}$ and $y<x$, use the converse half of the same
        local rule.
  \item If $\forall\DR.D\in T_{\sigma(x)}$ and $x\perp y$, choose the originating $\DR$
        birth cut. Both endpoint types lie in the inclusive cones of its control states.
        The cross-cone $\DR$ compatibility condition puts $D$ in $T_{\sigma(y)}$.
  \item $\EQ$ is local.
  \item The $\PO$ case is impossible by the fragment predicate.
\end{itemize}

\subsection{Existential restrictions}

For every non-$\EQ$ demand $\exists r.D$ at $x$, the chosen strategy creates a child $y$
whose type contains $D$. The geometry proves that $R(x,y)=r$: this is immediate for
$\PP$ and $\PPI$, the seed definition for $\DR$, and protected overlap for $\PO$.
An $\EQ$ demand is realized by $x$ itself.

\subsection{One-way truth lemma}

Interpret each concept name $P$ by
\[
  P^{M}=\{x\mid P\in T_{\sigma(x)}\}.
\]
Induction on $D\in\Cl$ proves
\[
  D\in T_{\sigma(x)}\Longrightarrow M,x\models D.
\]
Only this direction is required. Boolean cases follow from Hintikka consistency;
existential and universal cases use the preceding lemmas. If the accepted root signature
contains the prepared index of $\mathsf{nnf}(C_0)$, the root satisfies the normalized
concept; \texttt{sat\_nnf\_iff} transfers this back to $C_0$.

\textbf{Soundness theorem.}
\[
  \mathsf{ForallPOFree}(C_0)\wedge
  \mathsf{coneSatCoreB}(C_0)=\mathsf{true}
  \Longrightarrow \exists M,x,\;M,x\models C_0.
\]

\section{Completeness: every model yields surviving signatures}

Let $M$ be any model and $x$ a point. Define its realized support type and lower spectrum:
\[
  T_x=\{D\in\Cl\mid M,x\models D\},
  \qquad
  S_x=\{T_y\mid y\PP x\}.
\]
Since $\Ty$ is finite, $S_x$ is a finite subset of $\Ty$ even when the model has infinitely
many points. Put $q_x=(T_x,S_x)$.

\subsection{Static admissibility}

The truth definition gives local Hintikka consistency and the $\EQ$ conditions. If
$U=T_y\in S_x$, then $y\PP x$. Every $\forall\PP$ obligation at $y$ holds at $x$, and
every $\forall\PPI$ obligation at $x$ holds at $y$. Thus $q_x$ is statically admissible.

\subsection{Source transitions}

Suppose $y$ witnesses an existential demand at $x$.

\begin{itemize}
  \item If $x\PP y$, then every $z\leq x$ is a strict lower point of $y$, including $x$.
        Therefore $\cone(q_x)\subseteq S_y$.
  \item If $x\PPI y$, the dual inclusion $\cone(q_y)\subseteq S_x$ holds.
  \item If $x\DR y$, then every $z\leq x$ and $w\leq y$ is disjoint. This follows from
        the RCC5 singleton cells
        \[
          \operatorname{comp}(\PP,\DR)=\{\DR\},
          \qquad
          \operatorname{comp}(\DR,\PPI)=\{\DR\}.
        \]
        Hence every $\DR$ universal body transfers across all pairs in
        $\cone(q_x)\times\cone(q_y)$.
  \item If $x\PO y$, the target body holds at $y$ and no universal $\PO$ transfer exists.
  \item If $x\EQ y$, then $x=y$ and the demand is local.
\end{itemize}

Thus every source witness induces a compatible transition between realized signatures.

\subsection{Post-fixed-point argument}

Let
\[
  R_M=\{q_x\mid x\in M\}\subseteq\mathsf{Static}.
\]
In Lean, $R_M$ is a proof-level \texttt{Set ConeState}, defined by an existential source
point; it is not a computable \texttt{Finset} over an arbitrary domain.
Every existential demand at every state in $R_M$ has a target again in $R_M$, so
\[
  R_M\subseteq\Phiop(R_M).
\]
For a monotone operator on a finite powerset, every post-fixed set lies below the greatest
fixed point. Therefore $R_M\subseteq X_\infty$. If $M,x_0\models C_0$, semantic NNF
correctness gives $\mathsf{nnf}(C_0)\in T_{x_0}$, so an accepting survivor exists.

\textbf{Completeness theorem.}
\[
  \mathsf{ForallPOFree}(C_0)\wedge
  \bigl(\exists M,x,\;M,x\models C_0\bigr)
  \Longrightarrow \mathsf{coneSatCoreB}(C_0)=\mathsf{true}.
\]

This proof uses neither a finite model property nor a patchwork theorem. It compresses
only the \emph{local source signature}, not the source domain or its witnesses.

\section{Algorithm and conservative bound}

With $N\leq 2^m$ and $K=N2^N$, enumerate all $K$ raw signatures, filter static
admissibility, precompute compatible transition triples, and iterate elimination for at
most $K$ rounds.

A deliberately conservative implementation bound is:
\[
  O(K^2N^2m)\quad\text{for compatibility precomputation,}
\]
\[
  O(K^3m)\quad\text{for naive elimination.}
\]
Since $K\leq 2^{m+2^m}$, this gives a deterministic
\[
  2^{2^{O(m)}}
\]
time upper bound. This is an upper bound for the proposed implementation, not a claimed
tight complexity classification. The release theorem should state only decidability
until a Lean cost model and a separately reviewed complexity proof are available.

Several engineering improvements are safe but should follow correctness: bit-vector
types, indexed demand buckets, reverse dependency lists, and work-list elimination.

\section{Lean architecture and theorem graph}\label{sec:lean}

\subsection{Modules}

\begin{longtable}{>{\ttfamily\small}p{0.42\linewidth}p{0.52\linewidth}}
\toprule
\textnormal{\textbf{Module}} & \textbf{Responsibility} \\
\midrule
\endfirsthead
\toprule
\textnormal{\textbf{Module}} & \textbf{Responsibility} \\
\midrule
\endhead
RCC5/Atom.lean & atoms, converse, finite enumeration \\
RCC5/CompTable.lean & composition table and singleton-cell lemmas \\
RCC5/Network.lean & abstract strong RCC5 networks \\
RCC5/OrderedDisjoint.lean & normal form and composition closure \\
RCC5/TokenRepresentation.lean & exact region-set representation \\
Logic/Syntax.lean & raw concepts and role restrictions \\
Logic/NNF.lean & semantic NNF correctness and inverse-role normalization \\
Logic/Prepared.lean & normalized root code and positive closure package \\
Logic/PositiveClosure.lean & positive subformula closure and finiteness \\
Logic/POFree.lean & syntactic predicate and reflection \\
Logic/Semantics.lean & abstract and concrete model satisfaction \\
Logic/SupportType.lean & Hintikka types, enumeration, reflection \\
Logic/Demand.lean & finite existential demands \\
ConeScheme/State.lean & signatures, spectra, inclusive cones \\
ConeScheme/StateEnum.lean & exhaustive finite enumeration \\
ConeScheme/LocalRules.lean & static admissibility and Boolean checker \\
ConeScheme/Transition.lean & five role-specific compatibility rules \\
ConeScheme/Scheme.lean & abstract valid scheme with proof-independent targets \\
ConeScheme/Occurrence.lean & fresh demand-word domain \\
ConeScheme/VerticalComponent.lean & component tags and bridge uniqueness \\
ConeScheme/BaseOrder.lean & height, transitive closure, strict order \\
ConeScheme/DisjointClosure.lean & $\DR$ cuts and downward closure \\
ConeScheme/ProtectedPO.lean & exact residual overlap at $\PO$ births \\
ConeScheme/Frame.lean & ordered-disjoint and RCC5 frames \\
ConeScheme/Model.lean & valuation from support labels \\
ConeScheme/SpectrumInvariant.lean & actual lower types lie in spectra \\
ConeScheme/TruthLemma.lean & one-way semantic induction \\
ConeScheme/SchemeSoundness.lean & every valid abstract scheme has a model \\
ConeScheme/FiniteGfp.lean & generic reductive finite fixed-point facts \\
ConeScheme/Elimination.lean & reflected viability and survivor computation \\
ConeScheme/StrategyExtraction.lean & accepting survivors yield an abstract scheme \\
ConeScheme/SourceSignature.lean & signatures extracted from arbitrary models \\
ConeScheme/SourceTransitions.lean & source witness compatibility \\
ConeScheme/Soundness.lean & accepted strategy unfolds to a model \\
ConeScheme/Completeness.lean & realized signatures survive \\
ConeScheme/Decision.lean & public Boolean correctness and decidability \\
Regression/*.lean & adversarial instances and mutation guards \\
Audit/PublicTheorems.lean & only supported public exports \\
Audit/AxiomReport.lean & axiom checks for exported theorems \\
Audit/ExecutableDeps.lean & executable import and code-generation checks \\
\bottomrule
\end{longtable}

The former multi-tier extraction should move below \texttt{Legacy/MultiTier/}. The public
decision theorem must not import anything from \texttt{Legacy}.

\subsection{Core declarations}

\begin{lstlisting}
structure Prepared (C : Concept) where
  nnfC     : NNFConcept
  closure  : Finset NNFConcept
  rootIx   : ClosureIx closure
  nnf_eq   : nnfC = nnf C
  cl_eq    : closure = positiveClosure nnfC
  root_eq  : decode rootIx = nnfC

abbrev TypeCode (p : Prepared C) := Finset (ClosureIx p.closure)

structure ConeState (p : Prepared C) where
  ty       : TypeCode p
  lower    : Finset (TypeCode p)

def ConeState.cone (q : ConeState p) : Finset (TypeCode p) :=
  insert q.ty q.lower

abbrev DemandCode (p : Prepared C) := ClosureIx p.closure
def LocalOK (q : ConeState p) : Prop := ...
def StepOK (q : ConeState p) (d : DemandCode p)
    (q' : ConeState p) : Prop := ...

structure ConeScheme (p : Prepared C) where
  live     : Finset (ConeState p)
  root     : ConeState p
  root_mem : root ∈ live
  root_has : p.rootIx ∈ root.ty
  next     : ConeState p -> DemandCode p -> ConeState p
  next_mem : ...
  next_ok  : ...

def prune (X : Finset (ConeState p)) :=
  X.filter fun q => viableInB p X q
def survivors (p : Prepared C) : Finset (ConeState p) := ...

theorem regular_forest_model_iff
    (p : Prepared C) (hfree : ForallPOFree C) :
    (exists q, q ∈ survivors p /\ p.rootIx ∈ q.ty) <->
      Satisfiable C

theorem satisfiableDecidable (hfree : ForallPOFree C) :
    Decidable (Satisfiable C) :=
  if h : coneSatCoreB C = true then
    isTrue ((coneSatCoreB_correct hfree).mp h)
  else
    isFalse (fun hs => h ((coneSatCoreB_correct hfree).mpr hs))
\end{lstlisting}

\subsection{Named proof obligations}

The critical theorem spine is:

\begin{enumerate}
  \item \texttt{sat\_nnf\_iff},
        \texttt{inverseRole\_nnf\_conv}, and normalized-root preparation;
  \item Boolean reflection:
        \begin{itemize}
          \item \texttt{forallPOFreeB\_reflect};
          \item \texttt{supportTypeOKB\_reflect};
          \item \texttt{stateLocalOKB\_reflect};
          \item \texttt{demandsB\_complete};
          \item \texttt{transitionOKB\_reflect};
          \item \texttt{hasTargetB\_reflect};
          \item \texttt{viableInB\_reflect};
          \item \texttt{prune\_mem\_iff};
          \item \texttt{acceptingB\_reflect}.
        \end{itemize}
  \item \texttt{survivors\_fixed},
        \texttt{postfixed\_subset\_survivors}, and
        \texttt{survivors\_to\_scheme};
  \item \texttt{child\_unique\_parent}, \texttt{baseOrder\_raises\_height}, and
        \texttt{unfold\_lt\_rank};
  \item Geometry and representation:
        \begin{itemize}
          \item \texttt{dr\_cut\_separates};
          \item \texttt{occurrence\_orderedDisjoint};
          \item \texttt{orderedDisjoint\_rcc5};
          \item \texttt{tokenRepresentation\_exact}.
        \end{itemize}
  \item \texttt{protectedPO\_residual};
  \item \texttt{actualLower\_mem\_spectrum};
  \item \texttt{truth\_lemma} and \texttt{scheme\_sound}, both carrying the
        normalized-closure fragment hypothesis;
  \item \texttt{sourceSig\_localOK}, the five
        \texttt{source\_role\_step\_ok} lemmas, and
        \texttt{sourceSig\_survives};
  \item \texttt{regular\_forest\_model\_iff},
        \texttt{coneSatCoreB\_correct}, and \texttt{satisfiableDecidable}.
\end{enumerate}

The bundled JSON ledger records each obligation with dependencies, evidence status, and
its release gate. CI should reject a theorem marked complete when a dependency is not
complete.

\section{Regression suite}

The proof should be developed against semantic regressions, not only positive examples.

\begin{center}
\renewcommand{\arraystretch}{1.25}
\begin{tabularx}{\linewidth}{>{\bfseries}p{0.23\linewidth}XX}
\toprule
Regression & Required behavior & Mutation caught \\
\midrule
No finite model & Accept $C_\infty$ via recurrent control and fresh occurrences. &
Identifying recurrent states with points. \\
Mixed vertical cycle & Accept recurrent mixed control patterns while proving the unfolded
occurrence order acyclic. &
Identifying control states or reversing $\PPI$ orientation. \\
Protected PO & A $\PO$ demand is exactly residual at its birth bridge. &
Labelling $\PO$ before closure. \\
DR downward & Every pair below opposite $\DR$ sides remains $\DR$. &
Checking only seed endpoints. \\
Role-aware demand & Identical target concepts under different roles remain distinct. &
Indexing witnesses only by concept body. \\
D1 multi-group cycle & Replays the known borrowed-edge cycle attack. &
Reintroducing shared borrowed vertical edges. \\
D2 seven-point & Replays the incompatible blocked-source/target instance. &
Identifying target occurrences with old points. \\
E anchor compatibility & Replays the crossed-anchor incompatibility. &
Adding global anchors or shared horizontal witnesses. \\
\bottomrule
\end{tabularx}
\end{center}

For every successful theorem, add at least one mutation test that deletes or weakens its
decisive hypothesis. In particular, the $\DR$ completeness proof must fail if either
singleton composition cell is removed, and protected overlap must fail if component
bridge uniqueness is removed.

\section{Milestones and go/no-go gates}

\begin{center}
\renewcommand{\arraystretch}{1.25}
\begin{tabularx}{\linewidth}{>{\bfseries}p{0.10\linewidth}p{0.18\linewidth}Xp{0.13\linewidth}}
\toprule
Gate & Focus & Exit criterion & Estimate \\
\midrule
G0 & Source recovery & Existing syntax, semantics, RCC5 table, and theorem names compile. &
1--2 weeks \\
G1 & Finite control & Enumeration, reflected checks, GFP, and certificate extraction compile. &
2--3 weeks \\
G2 & Occurrence geometry & Strict order, $\DR$ separator, protected $\PO$, and normal form compile. &
3--5 weeks \\
G3 & Soundness & Spectrum invariant and truth lemma compile with no admitted facts. &
3--4 weeks \\
G4 & Completeness & Source signatures are a post-fixed point for every model. &
2--4 weeks \\
G5 & Public theorem & Boolean correctness and \texttt{Decidable} instance compile. &
2--3 weeks \\
G6 & Audit and release & Dual clean builds, regressions, axiom report, and package manifest pass. &
2--3 weeks \\
\bottomrule
\end{tabularx}
\end{center}

A realistic estimate is four to six focused months after source recovery. G2 is the first
hard go/no-go gate. Do not invest in optimized enumeration or complexity certification
until \texttt{protectedPO\_residual} and \texttt{occurrence\_orderedDisjoint} compile.
If G2 fails, geometry is the first suspected fault. A counterexample may nevertheless
show that the finite signature or transition predicate is too weak and require a redesign
of G1.

\section{Certification and release gates}\label{sec:audit}

\subsection{Logical hygiene}

The release build must:

\begin{itemize}
  \item contain no \texttt{sorry}, \texttt{admit}, \texttt{sorryAx}, or new axioms;
  \item contain no unsafe declaration in the proof dependency cone;
  \item compare \texttt{\#print axioms} for every public theorem against an exact
        allowlist, and separately scan all declarations for unused custom axioms;
  \item keep classical choice explicit and outside the Boolean decider;
  \item use kernel \texttt{by decide} for release regressions, or explicitly account for
        any extra trust introduced by \texttt{native\_decide};
  \item use warnings as errors;
  \item prove Boolean reflection for all finite predicates used by the executable.
\end{itemize}

\subsection{Dependency hygiene}

Build the project twice: once with the full tree and once from a minimal public import
that excludes \texttt{Legacy}. Generate the transitive import graph and fail CI if the
executable finite-control modules reach \texttt{SourceSignature}, \texttt{Completeness},
or \texttt{Legacy}. Run a generated-code evaluation test for \texttt{coneCheck}. Pin the
Lean toolchain and all Lake dependencies in the package manifest.

\subsection{Reproducibility}

The release archive should include:

\begin{itemize}
  \item exact Lean and Lake versions;
  \item a clean-build script;
  \item the theorem/axiom report;
  \item all regression files;
  \item a SHA-256 manifest;
  \item this report and the machine-readable obligation ledger.
\end{itemize}

Run the clean build on two independent runners. Extract the archive into a new directory
before the final build, so no untracked local file can mask a missing dependency.

\section{Risk register}

\begin{longtable}{>{\bfseries}p{0.24\linewidth}p{0.47\linewidth}p{0.20\linewidth}}
\toprule
Risk & Mitigation & Decision point \\
\midrule
Positive closure omits a needed body &
Prove structural closure lemmas before defining support types; add a theorem that every
semantic induction child belongs to $\Cl$. & G1 \\
Spectrum transitivity is under-specified &
State the invariant with inclusive cones and prove one-edge propagation before
transitive closure. & G3 \\
$\PPI$ height orientation is reversed &
Centralize role orientation in one definition and test mixed $\PP/\PPI$ words. & G2 \\
$\DR$ closure becomes reflexive &
Prove the cut separator before defining the ordered-disjoint instance. & G2 \\
$\PO$ birth is later ordered or disjoint &
Use component tags and unique bridges; prove protected overlap as a standalone theorem. &
G2 \\
Abstract and concrete semantics diverge &
Complete the ordered-token representation and prove atom-by-atom equivalence. & G2/G3 \\
Classical proof is mistaken for executable code &
Require reflected Booleans and inspect generated evaluator dependencies. & G1/G5 \\
Complexity claim outruns certification &
Publish decidability first; keep the double-exponential estimate explicitly conservative. &
G6 \\
\bottomrule
\end{longtable}

\section{Relation to the literature}

The plan uses standard ideas but the central combination is bespoke. RCC5 and its
composition calculus originate with Randell, Cui, and Cohn~\cite{RCC92}; formal model and
composition-table issues are analyzed by Li and Ying~\cite{LiYing2003}, while tractable
RCC5 sublanguages were classified by Jonsson and Drakengren~\cite{JD1997}. Type
elimination and tableau methods for description logics motivate the finite control layer,
including Donini and Massacci~\cite{DM2000} and the treatment of inverse and transitive
roles by Horrocks and Sattler~\cite{HS1999}. Regular grammar logics with converse provide
nearby automata and tableau methodology~\cite{DN2005,Nguyen2007}; Fischer--Ladner closure
and automata-theoretic modal decision procedures are further background
~\cite{FL1979,VW1986}. Lean 4 is described by de Moura and Ullrich~\cite{Lean4}.

None of these sources proves the theorem targeted here. In particular, the
fresh-occurrence cone scheme, the spectrum transition conditions, and the protected-$\PO$
regular-forest theorem require an original proof. The literature supports the surrounding
methods and RCC5 semantics, not the transfer step.

\section{Final recommendation}

Proceed with the cone-scheme architecture and freeze the previous multi-tier development
as legacy evidence. Implement G1 and G2 first. Their interfaces are narrow enough that
two proof engineers can work independently: one on finite reflected control, one on
fresh-occurrence geometry. Join them only at the spectrum invariant.

The project should be declared complete only when:

\begin{enumerate}
  \item the regular-forest model equivalence is proved with the exact fragment predicate;
  \item the executable Boolean decider is connected to that equivalence;
  \item the public theorem has an acceptable axiom report;
  \item every adversarial regression passes in a clean extracted archive.
\end{enumerate}

This is a materially different proof, not a patch to the blocked extraction. Its main
advantage is conceptual: the finite object records \emph{how to generate} a model, while
the infinite unfolding supplies the fresh geometry that RCC5 proper-part chains require.

\appendix

\section{Proof-obligation dependency summary}

\begin{center}
\renewcommand{\arraystretch}{1.25}
\begin{tabularx}{\linewidth}{>{\bfseries}p{0.18\linewidth}XX}
\toprule
Layer & Inputs & Output theorem \\
\midrule
Reflection & syntax, finite sets & exact Boolean predicates \\
Finite GFP & reflected local and transition checks & survivor fixed point \\
Geometry & chosen survivor strategy & ordered-disjoint occurrence frame \\
Soundness & geometry, spectrum invariant & accepted implies satisfiable \\
Source map & arbitrary model, singleton RCC5 cells & realized signatures post-fixed \\
Completeness & source map, greatest fixed point & satisfiable implies accepted \\
Decision & finite enumeration, equivalence & full decidability \\
\bottomrule
\end{tabularx}
\end{center}

The machine-readable file \texttt{blueprint/PROOF\_OBLIGATIONS.json} is authoritative for
individual names and dependencies. Status values in that file deliberately distinguish
paper argument, executable check, inherited result, and pending Lean proof.

\section{Package map}

\begin{itemize}
  \item \texttt{report/}: this source report and compiled PDF;
  \item \texttt{blueprint/DECISION\_PROCEDURE.md}: concise mathematical specification;
  \item \texttt{blueprint/LEAN\_MODULE\_MAP.md}: module ownership and import order;
  \item \texttt{blueprint/PROOF\_OBLIGATIONS.json}: theorem dependency ledger;
  \item \texttt{blueprint/check\_blueprint.py}: structural ledger validation;
  \item \texttt{regressions/}: finite adversarial probes from the round-two attack;
  \item \texttt{run\_checks.sh}: one-command blueprint and regression replay.
\end{itemize}

\section{Bibliography}

\begin{thebibliography}{10}

\bibitem{RCC92}
D. A. Randell, Z. Cui, and A. G. Cohn.
\newblock A spatial logic based on regions and connection.
\newblock In \emph{Principles of Knowledge Representation and Reasoning: KR'92},
pages 165--176, 1992.
\newblock \url{https://cdn.aaai.org/KR/1992/KR92-017.pdf}.

\bibitem{JD1997}
P. Jonsson and T. Drakengren.
\newblock A complete classification of tractability in RCC-5.
\newblock \emph{Journal of Artificial Intelligence Research}, 6:211--221, 1997.
\newblock \url{https://doi.org/10.1613/jair.385}.

\bibitem{LiYing2003}
S. Li and M. Ying.
\newblock Region Connection Calculus: its models and composition table.
\newblock \emph{Artificial Intelligence}, 145:121--146, 2003.
\newblock \url{https://doi.org/10.1016/S0004-3702(02)00372-7}.

\bibitem{DM2000}
F. M. Donini and F. Massacci.
\newblock EXPTIME tableaux for ALC.
\newblock \emph{Artificial Intelligence}, 124:87--138, 2000.
\newblock \url{https://doi.org/10.1016/S0004-3702(00)00070-9}.

\bibitem{HS1999}
I. Horrocks and U. Sattler.
\newblock A description logic with transitive and inverse roles and role hierarchies.
\newblock \emph{Journal of Logic and Computation}, 9(3):385--410, 1999.
\newblock \url{https://doi.org/10.1093/logcom/9.3.385}.

\bibitem{DN2005}
S. Demri and H. de Nivelle.
\newblock Deciding regular grammar logics with converse through first-order logic.
\newblock \emph{Journal of Logic, Language and Information}, 14:289--329, 2005.
\newblock \url{https://arxiv.org/abs/cs/0306117}.

\bibitem{Nguyen2007}
L. A. Nguyen.
\newblock A tableau calculus for regular grammar logics with converse.
\newblock In \emph{TABLEAUX 2007}, LNCS 4548, pages 421--436.
\newblock \url{https://doi.org/10.1007/978-3-540-73099-6_29}.

\bibitem{FL1979}
M. J. Fischer and R. E. Ladner.
\newblock Propositional dynamic logic of regular programs.
\newblock \emph{Journal of Computer and System Sciences}, 18:194--211, 1979.
\newblock \url{https://doi.org/10.1016/0022-0000(79)90046-1}.

\bibitem{VW1986}
M. Y. Vardi and P. Wolper.
\newblock Automata-theoretic techniques for modal logics of programs.
\newblock \emph{Journal of Computer and System Sciences}, 32:183--221, 1986.
\newblock \url{https://doi.org/10.1016/0022-0000(86)90026-7}.

\bibitem{Lean4}
L. de Moura and S. Ullrich.
\newblock The Lean 4 theorem prover and programming language.
\newblock In \emph{CADE 28}, 2021.
\newblock \url{https://leanprover.github.io/papers/lean4.pdf}.

\end{thebibliography}

\end{document}
